\begin{table}[tp]
\centering
\caption{Cross-country growth regressions, raw and purged inflation volatility}
\label{tab:crosscountry}
\begin{tabular}{lccccc}
\toprule
 & (1) & (2) & (3) & (4) & (5) \\
\midrule
Mean inflation $\bar\pi$ & $-0.0170^{***}$ &  &  & $-0.0114$ & $-0.0099^{*}$ \\
 & $(-3.52)$ &  &  & $(-1.28)$ & $(-1.65)$ \\
Raw $\sigma_\pi$ &  & $-0.0194^{***}$ &  & $-0.0082$ &  \\
 &  & $(-2.60)$ &  & $(-0.83)$ &  \\
Purged $\sigma_\pi/\sqrt{1+|\bar\pi|}$ &  &  & $-0.2032^{***}$ &  & $-0.1425^{**}$ \\
 &  &  & $(-3.28)$ &  & $(-2.18)$ \\
\midrule
Observations & 241 & 241 & 241 & 241 & 241 \\
Countries & 94 & 94 & 94 & 94 & 94 \\
$R^2$ & 0.602 & 0.600 & 0.604 & 0.603 & 0.608 \\
Period fixed effects & Yes & Yes & Yes & Yes & Yes \\
Country-clustered SE & Yes & Yes & Yes & Yes & Yes \\
Wild-cluster bootstrap $p$ &  &  &  & $0.43$ & $0.032$ \\
\bottomrule
\end{tabular}
\vspace{0.2em}
\begin{minipage}{0.90\linewidth}\footnotesize
\textit{Notes:} Each column is one growth regression on \citeauthor{barro1995inflation}'s panel
over his own periods---1965--75, 1975--85 and 1985--90; the dependent variable is the country's
average annual growth of real GDP per capita over the period, in 1985 international prices. Every
column includes his full control set, period fixed effects, and standard errors clustered by
country; $t$-statistics are in parentheses. The columns differ only in which summary of the
inflation experience enters: \textbf{(1)} mean inflation $\bar\pi$ alone; \textbf{(2)} the raw
within-period standard deviation $\sigma_\pi$ alone; \textbf{(3)} the purged measure alone;
\textbf{(4)} mean inflation together with the raw standard deviation; \textbf{(5)} mean inflation
together with the purged measure. Comparing~(4) with~(5) is the point of the exercise: the raw
standard deviation is virtually zero conditional on the level, as \citeauthor{barro1995inflation}
reports, and masks it; the purged measure is significant at five percent and the level's own
coefficient sharpens beside it.
\emph{Raw} is $\sigma_\pi$; \emph{purged} divides it by $\sqrt{1+|\bar\pi|}$. That deflator uses no
inflation target---no country in these periods had announced one---and contains nothing estimated:
it is $\sqrt{1+|\bar\pi|}$ exactly, with $\bar\pi$ in percentage points, the units in which
the envelope of \eqref{eq:var_decomp} is fitted. The dependent variable is a fraction: in column~(5) a
one-standard-deviation rise in the purged measure is worth $0.28$ percentage points of
annual growth, against $0.15$ for the raw measure in column~(4). The wild-cluster row
reports restricted Rademacher bootstrap $p$-values ($999$ draws, clustered by country)
for the volatility term of columns~(4) and~(5).
The controls are \citeauthor{barro1995inflation}'s own: initial income, male and female school
attainment at the secondary and higher levels, the interaction of initial income with human
capital as he defines it, life expectancy, fertility, government consumption net of education and
defence, public education spending, the black-market premium, the terms-of-trade change, the
investment ratio, and the political-rights index with its square. The one control of his we cannot
match is the rule-of-law index, which is not in the \citeauthor{barro1995inflation}--Lee
distribution and which he himself uses among his instruments. Inflation moments are period means
and standard deviations of the annual log change in the consumer price index
\citep{muller2025global}, untrimmed, as in his own exercise. Script: \texttt{tab03\_cross\_country.py}.
$^{***}p<0.01$, $^{**}p<0.05$, $^{*}p<0.1$.
\end{minipage}
\end{table}
